% ----------------------
\subsection{Historical background and source information}
% ----------------------
	\label{JST-HB}
	
% -----------------------------------------------------------
\pagebreak{}
% -----------------------------------------------------------

\begin{landscape}

	% -----
	\subsubsection{Timeline of the New Translation (JSNTB 57--59)}
	% -----
		\label{JST-TIMELINE}
	
		\scriptsize
	
		\begin{longtable}{p{1in}p{1.5in}p{1in}p{0.75in}p{1.5in}p{3in}}
			
			% HEADER
			\textbf{\textsc{Date}} & \textbf{\textsc{Scripture Reference}} & \textbf{\textsc{Scribe}} & \textbf{\textsc{Location}} & \textbf{\textsc{Manuscript Reference}} & \textbf{\textsc{Comments}} \\ \hline
			
			% June 1830
			June 1830 & Moses 1 (Comm. of Christ D\&C 22) & Oliver Cowdery & Fayette, N.Y. or Harmony, Pa. & OT1-1, line 1 to OT1-3, line 14 & Original dictation. Heading reads ``A Revelation given to Joseph the Revelator June 1830.'' \\ \hline
			
			% Between June and 21 October 1830
			Between June and 21 October 1830 & Moses 2:1--5:43; Genesis 1:1--4:18; \textit{I.V.} Genesis 1:1--5:28 & Oliver Cowdery & Harmony, Pa. and Fayette, N.Y. & OT1-3, line 15 to OT1-10, line 5 & Original dictation. Chapter heading on OT1-3 reads ``A Revelation given to the Elders of the Church of Christ On the First Book of Moses.'' Chapter heading on OT1-8 reads ``Chapter 2---A Revelation concerning Adam after he had been driven out of the garden of Eden.'' \\ \hline
			
			% 21 October 1830
			21 October 1830 & Moses 5:43--51; Genesis 4:18--24; \textit{I.V.} Genesis 5:29--37 & John Whitmer & Fayette, N.Y. & OT1-10, lines 6 to 23 & Original dictation. Probably translated and written in one day. \\ \hline
			
			% 30 November 1830
			30 November 1830 & Moses 5:52--6:18; Genesis 4:25--5:11; \textit{I.V.} Genesis 5:38--6:16 & John Whitmer & Fayette, N.Y. & OT1-10, line 24 to OT1-11, line 40 & Original dictation. Probably translated and written in one day. \\ \hline
			
			% 1 December 1830
			1 December 1830 & Moses 6:19--52; Genesis 5:12--21; \textit{I.V.} Genesis 6:17--53 & Emma Smith & Fayette, N.Y. & OT1-11, line 41 to OT1-14, line 1 & Original dictation. Probably translated and written in one day. Emma Smith commanded, ``And thou shalt go with him [Joseph Smith]...and be unto him for a scribe'' (\hyperref[DC25-6]{DC 25:6}). \\ \hline
			
			% Between 1 and 10 December 1830
			Between 1 and 10 December 1830 & Moses 6:52--7:1; \textit{I.V.} Genesis 6:53--7:1 & John Whitmer & Fayette, N.Y. & OT1-14, line 1 to OT1-15, line 16 & Original dictation. Heading on OT1-14 reads ``The Plan of Salvation.'' \\ \hline
			
			% Ca. 10 December 1830 to 7 March 1831
			Ca. 10 December 1830 to 7 March 1831 & Moses 7:2--8:30; Genesis 5:22--24:41; \textit{I.V.} Genesis 7:2--24:42 & Sidney Rigdon & Fayette, N.Y. and Kirtland, Ohio & OT1-15, line 16 to OT1-61, line 5 & Original dictation. Sidney Rigdon commanded to ``write for him [Joseph Smith]; and the scriptures shall be given, even as they are in mine own bosom'' (\hyperref[DC35-20]{DC 35:20}). \\ \hline
			
			% 8 March--7 April 1831
			8 March--7 April 1831 & Matthew 1:1--9:1 & Sidney Ridgon & Kirtland, Ohio & NT1-1, line 1 to NT1-21, line 2 & Original dictation. Heading reads ``A Translation of the New Testament translated by the power of God.'' Joseph Smith commanded 7 March 1831 to begin translation of New Testament (\hyperref[DC45-60]{DC 45:60--61}). \\ \hline
			
			% Ca. 8 March--5 April 1831
			Ca. 8 March--5 April 1831 & Moses 1--8; Genesis 1:1--24:41; Comm. of Christ D\&C 22; \textit{I.V.} Genesis 1:1--24:42 & John Whitmer (copyist) & Kirtland, Ohio & OT2-1, line 1 to OT2-59, line 1 & Transcribed from OT1-1, line 1 to OT1-61, line 5. John Whitmer commanded on 8 March 1831 to ``assist you, my servant Joseph, in transcribing all things which shall be given you'' (\hyperref[DC47-1]{DC 47:1}). \\ \hline
			
			% Between 4 and 7 April 1831
			Between 4 and 7 April 1831 & Matthew 1:1--9:1 & John Whitmer (copyist) & Kirtland, Ohio & NT2.1-1, line 1 to NT2.1-16, line 16 & Transcribed from NT1-1, line 1 to NT1-21, line 2. \\ \hline
			
			% 7 April--before 19 June 1831
			7 April--before 19 June 1831 & Matthew 9:2--26:71 & Sidney Rigdon & Kirtland, Ohio & NT1-21, line 3 to NT1-63, line 12 & Original dictation. Probably completed prior to Joseph Smith and other Church leaders departing to Zion (Jackson County, Missouri) on 19 June 1831. \\ \hline
			
			% Between 7 April and 26 September 1831
			Between 7 April and 26 September 1831 & Matthew 9:2--26:1 & John Whitmer (copyist) & Kirtland, Ohio & NT2.1-16, line 15 to NT2.1-49, line 34 & Transcribed from NT1-21, line 3 to NT1-59, line 39. \\ \hline
			
			% 26 September--before 20 November 1831
			26 September--before 20 November 1831 & Matthew 26:1--Mark 9:1; \textit{I.V.} Matthew 26:1--Mark 8:44 & John Whitmer & Hiram, Ohio & NT2.2-1, line 1 to NT2.2-24, line 23 & Original dictation. Written prior to John Whitmer departing Ohio (with Oliver Cowdery) for Jackson County, Missouri, on 20 November 1831. \\ \hline
			
			% Ca. 20 November 1831 to 16 February 1832
			Ca. 20 November 1831 to 16 February 1832 & Mark 9:2--John 5:29; \textit{I.V.} Mark 9:1--John 5:29 & Sidney Rigdon & Hiram, Ohio & NT2.2-24, line 24 to NT2.4-114, line 18 & Original dictation. Joseph Smith and Sidney Rigdon received the Vision (\hyperref[DC76]{DC 76}) on 16 February 1832 while translating John 5:29. \\ \hline
			
			% Between 16 February and 24 March 1832
			Between 16 February and 24 March 1832 & John 5:30--Revelation 11:4 & Sidney Rigdon and Scribe A & Hiram, Ohio & NT2.4-114, line 18 to NT2.4-152, line 2 & Original dictation. Short notation system begun with John 6 at NT2.4-115, line 10. Scribe A wrote on NT2.4-136, 139--42, 147--149. \\ \hline
			
			% Between 20 and 31 July 1832
			Between 20 and 31 July 1832 & Revelation 12:1--22:9 & Frederick G. Williams & Hiram, Ohio & NT2.4-152, line 3 to NT2.4-154, line 11 & Original dictation. Frederick G. Williams began as Joseph Smith's scribe on 20 July 1832. On 31 July 1832, Joseph Smith wrote: ``We have finished the translation of the New Testament'' (see \hyperref[EMS-1832-JS-31-JUL-1832]{this letter}). \\ \hline
			
			% Between 20--31 July 1832 and 2 February 1833
			Between 20--31 July 1832 and 2 February 1833 & Matthew--Revelation & Sidney Rigdon and Frederick G. Williams & Hiram and Kirtland, Ohio & Review of NT2.1 to NT2.4 & Additional revisions to New Testament text. On 2 February 1833, Frederick G. Williams noted: ``This day completed the translation and the reviewing of the New Testament'' (KCMB 8). \\ \hline
			
			% 20--31 July 1832--2 July 1833
			20--31 July 1832--2 July 1833 & Genesis 24:41--Nehemiah 10:30; \textit{I.V.} Genesis 24:42--Nehemiah 10:30 & Frederick G. Williams & Hiram and Kirtland, Ohio & OT2-59, line 1 to OT2-81, line 20 & Original dictation. On 31 July 1832, after announcing completion of New Testament, Joseph Smith wrote: ``We are making rapid strides in the old book [Old Testament]'' (see \hyperref[EMS-1832-JS-31-JUL-1832]{this letter}). Short notation system begun with Genesis 25:7 at OT2-60, line 19. \\ \hline
			
			% Between late July 1832 and 2 July 1833
			Between late July 1832 and 2 July 1833 & Nehemiah 11--Psalm 10 & Joseph Smith & Kirtland, Ohio & OT2-81, line 21 to OT2-83, line 32 & Original dictation. The Prophet acted as his own scribe. \\ \hline
			
			% Between late July 1832 and 2 July 1833
			Between late July 1832 and 2 July 1833 & Psalm 11--Malachi & Frederick G. Williams & Kirtland, Ohio & OT2-83, line 33 to OT2-119, line 5 & Original dictation. Joseph Smith wrote for himself on OT2-86, lines 9 to 21 (Psalm 16), and Sidney Rigdon was scribe for a few lines on OT2-111--12 (Jeremiah 18). At the end of Malachi, scribe Frederick G. Williams wrote: ``Finished on the 2d day of July 1833.'' \\
			
		\end{longtable}
	
\end{landscape}\clearpage\normalsize

% ----------------------
\subsection{Old Testament Revision 1}
% ----------------------
	\label{JST-OTR1}
	
	\label{JST-OTR1}

% -----
\subsubsection{OTR1 Genesis}
% -----
	\label{JST-OTR1-GENESIS}
	
	% --
	\paragraph{OTR1 Moses 1}
		% STATUS
			% this chapter has been copied/pasted in its entirety and I have finished doing all the markup. It's done.
	% --
		\label{JST-OTR1-MOSES-1}
		
		A Revelation given to Joseph the Revelator June 1830
		
		The words of God which he \sout{gave} <spake> unto Moses at a time when Moses was caught up into an exceeding high Mountain \& he saw God face to face \& he talked with him \& the glory of God was upon Moses therefore Moses could endure his presence \& God spake unto Moses saying Behold I I am the Lord God Almighty \& endless is my name for I am without beginning of days or end of years \& is this not endless \& behold thou art my Son Wherefore look \& I will shew thee the workmanship of mine hands but not all for my works are without end \& also my wo[r]ds for they never cease wherefore no man can behold all my works except he behold <all> my glory \& no man can behold all my glory \& afterwards remain in the flesh \& I have a work for thee Moses my Son \& thou art in similitude to \sout{my} <mine> only begotten \& mine only begotten is \& shall be for he is full of grace \& truth bu[t] there is none other God beside me \& all things are present with me for I know them all And now behold this one thing I shew unto thee Moses my son for thou art in the world \& now I shew it thee And it came to pass that Moses looked \& beheld the world upon which he was created \& [as] Moses beheld the world \& the ends thereof \& all the Children of men which was \& which was created of the same he greatly marvelled \& wondered \& the presence of God withdrew from Moses that his glory was not upon Moses \& Moses was left unto himself \& as he was left unto himself he fell unto the Earth, And it came to pass, that it was for the space of many hours before Moses did again receive his natural strength [lik]e unto man, \& he saith unto himself \sout{no} Now for this once I know that [m]an is nothing, which thing I never had supposed, but now mine eyes, mine own eyes but not mine eyes for mine eyes could not have beheld for I should have withered \& died in his presence but his glory was upon me \& I beheld his face for I was transfigered before him And now it came to pass that when Moses had said these words, behold, Satan came tempting him saying, Moses, Son of man, worship me, And it came to pass that Moses looked upon Satan \& saith, Who art thou for behold, I am a Son of God in the similitude of his only begotten, \& where is thy glory that I should worship thee, for, behold, I could not look upon God except his glory should come upon me, \& I were transfigered before him but I can look upon thee in the natural man, if not so surely blessed be the name of my God for his Spirit hath not altogether withdrawn from m[e] or else where is thy glory for it is blackness unto me \& I can Judge between thee \& God for God said unto me Worship God for him only shalt thou serve Get thee hence Satan deceive me not for God said unto me Thou art after the similitude of mine only begotten \& he also gave unto me commandment when he called unto me out of the burning bush Saying [c]all upon God in the name of mine only begotten \& worship me And again Moses saith I will not cease to call upon God I have other things to inquire of him for his glory has been upon me \& it is glory unto me wherefore I can judge betwixt him \& thee depart hence Satan And now when Mose[s] had said these words Satan cried with a loud [voice \&] wrent upon the Ear[th] \& comma[n]ded saying I am the only begotten worship me And it came to pass that Moses began to fear exceedingly \& as he began to fear he saw the bitterness of Hell Nevertheless calling upon God he received strength \& he commanded saying Depart hence Satan for this one God only will I worship which is the God of glory \& now Satan began to tremble \& the Earth shook \& Moses receiving strength called upon God saying In the name of Jesus Christ depart hence Satan And it came to pass that satan cried with a loud voice with weeping \& wailing \& gnashing of teeth \& departed hence yea from the presence of Moses that he beheld him not And now of this thing Moses bore record but because of wickedness it is not had among the children of men And it came to pass that when Satan had departed from the presence of Moses he lifted up his eyes unto Heaven being filled with the Holy Ghost which beareth record of the Father \& the Son \& calling upon the name of God he beheld again his glory for it was upon him \& he heard a voice saying Blessed art thou Moses for I the Almighty have chosen thee \& thou shalt be made stronger than the many waters for they shall obey thy command even as if thou wert God \& lo I am with you even to the end of thy days for thou shalt deliver my people from bondage even Israel my chose[n] And it came to pass as the voice was still speaking he cast his eyes \&beheld the Earth yea even all all the face of it \& there was not a particl[e] of it which he did not behold diserning it by the Spirit of God \& he beheld also the inhabitants thereof \& there was not a soul which he be[held] not \& he discerned them by the Spirit of God \& their numbers were grea[t] even as numberless as the sand upon the sea shore \& he beheld many lands \& each land was called Earth \& there were inhabitants upon the face there of And it came to pass that Moses called upon God saying tell me I pray thee why these things are so \& by what thou madest them \& behold the glory of God was upon Moses that Moses stood in the presence of God \& he talk[ed] with him face to face \& the Lord God said unto Moses For mine own purpose have I made these things here is wisdom \& it remaineth in me \& by the word of my power have I created them which is mine only begotten Son full of grace \& truth \& worlds without number have I created \& I also created them for mine own purpose \& by the same I created them which is mine only begotten \& the first man of all men have I called Adam which is many but only an account of this Earth \& the inhabitants thereof give I unto you for behold there are many worlds which have passed away by the word of my power \& there are many also which now stand \& numberless are they unto man bu[t] all things are numbered unto me for they are mine \& I know them And it came to pass that Moses spake unto the Lord saying Be merciful [u]nto thy servant O God \& tell me concerning this Earth \& the inhabitan[ts] [th]ereof \& also the H[eavens] \& then thy servant will be content
		
		And the Lord God spake unto Moses saying The Heavens there are many \& they cannot be numbered unto man but they are numbered unto me for they are mine \& as one Earth shall pass away \& the Heavens thereof even so shall another come And there is no end to my works neither my words for behold this is my work to my glory to the immortality \& the eternal life of man And now Moses my Son I will speak unto you concerning this Earth upon which thou standest \& thou shalt write the things which I shall speak \& in a day when the children of men shall esteem my words as naught \& take many of them from the Book which thou shalt write behold I will raise up another like unto \sout{theee} thee \& they shall be had again among the Children of men among even as many as shall believe These words was spoken unto Moses in the mount the name of which shall not be known among the Children of men And now they are also spoken unto you shew them not unto any except them that believe <Amen>
 
		A Revelation given to the Elders of the Church of Christ On the first Book of Moses \sout{given to Joseph the Seer}
		
	% --
	\paragraph{OTR1 Genesis 1}
		% STATUS
			% this chapter has been copied/pasted in its entirety and I have finished doing all the markup. It's done.
	% --
		\label{JST-OTR1-GENESIS-1}
		
		\begin{tightcenter}
		Chapter first
		\end{tightcenter}
		
		And it came to pass that the Lord spake unto Moses saying Behold I reveal unto you concerning this Heaven \& this Earth write the words which I speak I am the beginning \& the end the Almighty God by mine only begotten I created these things yea in the beginning I created the Heaven \& the Earth upon which thou standest \& the Earth was without form \& void \& I caused darkness to come up upon the face of the deep \& my Spirit moved upon the face of the waters for I am God \& I God said Let there be light \& there was light \& I God saw the light \& the light was good \& I God divided the light from the darkness \& I God called the light day \& the darkness I called night \& this I done by the Word of my power \& it was done as I spake \& the evening \& the morning were the first day And again I God said let there be a firmament in the \sout{midst} <midst> of the waters \& it was so even as I spake And I said Let it divide the waters from the waters \& it was done \& I God made the firmament \& divide[d] the waters yea the great waters under the firmament from the waters which were above the firmament \& it was so even as I spake And I God called the firmament Heaven \& the evening \& the morning were the second day And I God said Let the waters under the Heaven be gathered together unto one place \& it was so \& I God said Let there be dry land [\&] it was so \& I God called the dry land Earth \& the gathering together <of> the waters called I <the> Seas \& I God saw that all things which I had made were good \& I God said Let the Earth bring forth grass the herb yielding seed the fruit tree yielding fruit after his kind \& the tree yielding fruit whose seed should be in itself \sout{after} upon the Earth \& it was so even as I spake \& the Earth brought forth grass every herb yielding seed after his kind \& the tree yielding fruit whose seed should be in itself after it[s] kind--- $\diamond$
		
		And I God saw that all things which I had made were good \& the evening \& the morning were the third day And I God said Let there be lights in the firmament of the Heaven to divide the day from the night \& let them be for signs \& for seasons \& for days \& for years \& let them be for lights in the firmament of the Heaven to give light upon the Earth \& it was so \& I God made two great lights the greater light to rule the day \& the lesser light to rule the night \& the greater light was the Sun \& the lesser light was the moon \& the stars also was made even according to my word \& I God set them in the firmament of the Heaven to give light upon the Earth \& the Sun to rule over the day \& the moon to rule over the night \& to divide the light from the darkness \& I God saw that all things which I made were good \& the evening \& the morning were the fourth day \& I God said Let the waters bring forth abundantly the moveing creature that hath life \& fowl which may fly above the Earth in the open firmament of Heaven \& I God created great whales \& every living creature that moveth which the waters brought forth abundantly after their kind \& every winged fowl after his kind \& <I> God saw that all things which I had created were good \& I God blessed them saying be fruitful \& multiply \& fill the waters in the Seas \& let fowl multiply in the Earth \& the evening \& the morning were the fifth day \& I God said Let the Earth bring forth the living creature after his kind cattle \& creeping thing \& beast <of the eart[h]> after his kind \& it was so \& I God made the beast of the Earth after his kind \& cattle after their kind \& every thing which creepeth upon the Earth after his kind \& I God saw that all these things were good \& I God said unto mine only begotten which was with me from the beginning Let us make man in our image after our likeness \& it was so \& I God said Let them have dominion over the fish of the Sea \& over the fowl of the air \& over the cattle \& over all the Earth \& over every creeping thing that creepeth upon the Earth So I God created man in mine own image in the image of mine only begotten created I him male \& feemale created I them \& I God blessed them \& I God said unto them Be fruitful \& multiply \& replenish the Earth \& subdue it \& have dominion over the fish of the Sea \& over the fowl of the air \& over every living thing that moveth upon the Earth \& I God said unto man Behold I have given you every herb bearing seed which is upon the face of all the Earth \& every tree in the which shall be the fruit of a tree yielding seed to you it shall be for meat \& to every beast of the Earth \& to every fowl of the air \& to every thing that creepeth upon the Earth wherein I grant life there shall be given every clean herb for meat \& it was so even as I spake \& I God saw every thing that I had made \& behold all things which I had made were \sout{verry} <very> good \& the evening \& the morning were the sixth day------
		
	% --
	\paragraph{OTR1 Genesis 2}
		% STATUS
			% % this chapter has been copied/pasted in its entirety and I have finished doing all the markup. It's done.
	% --
		\label{JST-OTR1-GENESIS-2}
		
		Thus the Heaven \& the Earth were finished \& all the host of them \& on the seventh day I God ended my work \& all things which I had made \& I rested on the seventh day from all my work \& all things which I had made were finished \& I saw that they were good \& I God blessed the seventh day \& sanctified it because that in it I had rested from all my work which I God had created \& made And now behold I say unto you that these are the generations of the Heaven \& of the Earth when they were created in the day that I the Lord God made the Heaven \& the Earth \& every plant of the field before it was in the Earth \& every herb of the field before it grew for I the Lord God created all things of which I have spoken spiritually before they were naturally upon the face of the Earth for I the Lord God had not caused it to rain upon the face of the earth \& I the Lord God had created all the children of men \& not yet a man to till the ground for in Heaven created I them \& there was not yet flesh upon the Earth neither in the water neither in the air but I the Lord God spake \& there went up a mist from the Earth \& watered the whole face of the ground \& I the Lord God formed Man from the dust of the ground \& breathed into his nostrils the breath of life \& man became a living soul the first flesh upon the Earth the first man also nevertheless all things were before created but spiritually were they created [\&] made according to my word \& I the Lord God planted a garden eastward in Eden \& there I put the\sout{m} man whom I had formed \& out of the ground made I the Lord God to grow every tree naturally that is pleasant to the sight of man \& man could behold it \& they became also a living soul for it was \sout{spiritaull} spiritual in the day that I created it for it remaineth in the sphere which I created it yea even all things which I prepared for the use of man \& man saw that it was good for food \& I the Lord God placed the tree of life also in the midst of the garden \& also the tree of knowledge of good \& evil \& a River went out of Eden to water the garden \& from thence it was parted \& became into four heads \& I the Lord God called the name of the <first> Pison \& it compasseth the whole land of Havilah where there were created much gold \& the gold of that land was good \& there was bdellium \& the Onyx stone \& the name of the second River was called Gihon the same was it that compasseth the whole land of Ethiopia \& the name of the third River was Hiddekel that was it which goeth towards the east of Assyria \& the fourth River was Euphrates \& I the Lord God took the man \& put him into the garden of Eden \sout{\&} to dress it \& to keep it \& I the Lord God commanded the man saying Of every tree of the garden thou mayest freely eat but of the tree of the knowledge of good \& evil thou shalt not eat of it nevertheless thou mayest chose for thyself for it is given unto thee but remember that I forbid it for in the day that thou eatest thereof thou shalt surely die And I the Lord God said unto mine only begotten that it was not good that the man should be alone Wherefore I will mak[e] a help meet for him \& out of the ground I the Lord God formed every beast of the field \& every fowl of the air \& commanded that they should be brought unto Adam to see what he would call them \& they were also living souls \& it was breathed into them the breath of life \& whatsoever Adam called every living creature that was the name thereof \& Adam gave names to all cattle \& to the fowl of the air \& to every beast of the field but for Adam there was not found a help meet for him \& I the Lord God caused a deep sleep to fall upon Adam \& he slept \& I took one of his ribs \& closed up the flesh in the stead thereof \& the rib which I the Lord God had taken from man made I a woman \& brought her unto the man \& Adam said this I know now is bone of my bones \& flesh of my flesh she shall be called woman because she was taken out of man therefore shall a man leave his father \& his mother \& shall cleave unto his wife \& they shall be one flesh \& they were both naked the ma[n] \& his wife \& were not ashamed.
	
	% --
	\paragraph{OTR1 Genesis 14}
		% STATUS
			% this chapter has been copied/pasted in its entirety and I have finished doing all the markup. It's done.
	% --
		\label{JST-OTR1-GENESIS-14}
		
			\hypertarget{JST-OTR1-GENESIS-14-a}%
		And
			\marginnote{a}%
		it came to pass in the days of Amraphel king of Shinar and Arioch king of Elasar Chederlaomer king of Elam and Tidal king of Nations <that> these kings made war with Bera king of Sodom and with Bersha king of \sout{Sodom} Gomorrah Shinab king of Adma and Shemeber king of \sout{Zebioim} Zeboiim and the king of Bela which is Zoar all these were joined together in the vale of Siddim which is the Salt sea twelve years they Served Chederlaomer and in the thirteenth year they rebelled and in the fourteenth year came Chederlaomer and the kings that were with him and smote the Rephaims in Ashteroth Karnaim and <the> Zuzims in Ham <and> the \sout{Emmims} Emims in Shaveh Keriatheim and the Horites in <their> Mount Seir unto Elparan which was by the willderness and the<y> returned and came to Enmishpat which is Kadesh and smote all the country of the Amalakites and also the Amorites Hazezontamar and there went out the King of Sodom and the King of Gomorrah and the king of \sout{Amah} Admah and the King of Zeboiim and the King of Bela which is Zoar and the<y> joined battle with them in the Vale of Siddim with Chederlaomer King of Elam and with Tidal King of Nations and \sout{Armaphel} <Amraphel> King of Shinar, and \sout{Arioch} Arioch King of Ellasar four Kings with five and the vale of Siddim \sout{were} <was> filled with slime pits and the Kings of Sodom and Gomorrah fled and fell there and <they> that remained fled to the Mountain which was called \sout{Hanable} Hanabal and the<y> took all the goods of Sodom and Gomorrah and all their victuals and went their way
			\marginnote{b}%
		and
			\hypertarget{JST-OTR1-GENESIS-14-b}%
		the<y> took Lot Abram's brother's Son who dwelt in Sodom and his goods and departed and there came one that had escaped and told Abram the Hebrew the man of God for he dwelt <in> the plain of Mamree the Amorite brother of \sout{Escole and} Eshcol and brother of Ener and these were confederate with Abram and when Abram heard that Lot his brothers \sout{sons} Son was taken captive he armed his trained men and they which <were born> in his own house three Hundred and eighteen and persued unto Dan and he divided himself againsts them he and his men by night and smote them and persued them unto Hobah which was on the left hand of Demascus and he brought back \sout{all the goods} Lot his brothers son and all his goods and the women also and the people and the King of Sodom also went out to meet him after his return from the slaughter of Chederleomar and of the Kings that were with <him> at the valy of \sout{Shaveeh} Shaveh which was the Kings dale
			\marginnote{c}%
		and
			\hypertarget{JST-OTR1-GENESIS-14-c}%
		Melchizedeck King of Salam brought \sout{fourth} forth bread and wine and he break bread and blessed it and he blessed the wine he being the Priest of the most High God and he gave to Abram and he blessed him him and said blessed Abram thou art a man of the most \sout{high} High God \sout{possesor} possessor of heaven and earth and blessed is the name of the most High God which hath delivered \sout{thine} enemies into thine <thy> hand and \sout{he} Abram gave him tithes of all he had taken and the King of Sodom said unto Abram give me the persons and take the goods to thyself and Abram said to the King of Sodom I have lifted up my hand unto the Lord the most High God the possessor of heaven and earth and have sworn that I will not take \sout{from} of <thee> from a thread even to a shoe latchet and that I will not take any thing that is thine lest thou shouldst say I have made Abram rich save only that which the young men have eaten and the portion of the men which went with <me> Ener Eshcol <\sout{a} <\&>> Mamre let them take their portion. And Melchizedeck lifted up his voice and blessed Abram.
			\marginnote{d}%
		Now
			\hypertarget{JST-OTR1-GENESIS-14-d}%
		Melchizedeck was a man of faith who wrought righteousness and when a child he feared God and stoped the mouths of Lions and quenched the violence of fire%
			\footnote{%
				% \label{}%
				Speaking about people who exercised great faith, \hyperref[HEBREWS-11-33]{Hebrews 11:33--34} tell about those ``Who through faith subdued kingdoms, wrought righteousness, obtained promises, stopped the mouths of lions, Quenched the violence of fire, escaped the edge of the sword, out of weakness were made strong, waxed valiant in fight, turned to flight the armies of the aliens.'' 
				}
		and thus having been approved of God%
			\footnote{%
				% \label{}%
				Christ is also described as being ``a man approved of God'' in \hyperref[ACTS-2-22]{Acts 2:22}.
				}
		he was ordained \sout{<a>} a high Preist after the \hyperref[ORDER]{order} of the covenent which God made with Enoch it%
			\footnote{%
				% \label{}%
				I think the antecedent of this pronoun is ``order,'' rather than ``covenant.''
				}
		being after the \hyperref[ORDER]{order} of the Son of God which \hyperref[ORDER]{order} came not by man nor the will of men \sout{but of} neither by father nor Mother neither by begining of days nor end of years but of God and was delivered unto men by the calling of his own voice%
			\footnote{%
				% \label{}%
				Compare \hyperref[MOSES-5-58]{Moses 5:58} but especially \hyperref[MOSES-6-51]{Moses 6:51}. See also \hyperref[DC107-45]{DC 107:45}; 
				} 
		according to his own will unto as many as beleived on his name for God having sworn%
			\footnote{%
				% \label{}%
				Notice the ``having sworn,'' ``being ordained,'' ``having this faith,'' ``coming up unto,'' and other gerund constructions here---very similar to the many weird gerunds in \hyperref[ALMA-13-3]{Alma 13:3--9}.
				}
		unto Enoch and unto his seed with an oath by himself that every one being ordained after this \hyperref[ORDER]{order} and calling should have power by faith to break Mountains to divide the seas to dry up watters to turn them out of their course to put at defience the armies of nations to divide the earth to break every band to stand in the preasence of God to do all things according to his will according to his command subdue principalities and powers and this by the will of the Son of God which was from before the foundation of the world and men having this faith coming up unto this \hyperref[ORDER]{order} of <God> were translated and taken up into heaven.%
			\marginnote{e}%
		And
			\hypertarget{JST-OTR1-GENESIS-14-e}%
		now Melchizedeck was a Priest of this \hyperref[ORDER]{order} therefore he obtained peace in Salam and was called the prince of peace and his people wraught righteousness and obtained \sout{heaven} heaven and sought for the City of Enoch which God had before taken seperating it from the earth having reserved it unto the latter days or the end of the world and hath said and sworn with an oath that the heavens and the earth should come together and the sons of God should be tryed so as by fire and this Melchizedeck having thus established righteousness was called the King of heaven by his people or in other words the King of peace and he lifted up his voice and he blessed Abram being the high Priest and the keeper of the store house of God him whom God had appointed to receive tithes for the poor wherefore Abram paid unto him tithes of all that he had of all the riches which he posessed which God had given him more \sout{that} than that which he had need and it came to pass that God blessed Abram and gave unto him riches and honour and lands for an everlasting possession according to the covenent which he had made and according to the blessing wherewith Melchizedeck had blessed him.
	
% ----------------------
\subsection{Old Testament Revision, John Whitmer First Copy}
% ----------------------
	\label{JST-OTRJW}
	
	\input{gospel_scriptures_JST_OTRJW.tex}
	
% ----------------------
\subsection{Old Testament Revision 2}
% ----------------------
	\label{JST-OTR2}
	
	\input{gospel_scriptures_JST_OTR2.tex}
	
% ----------------------
\subsection{New Testament Revision 1}
% ----------------------
	\label{JST-NTR1}
	
	\input{gospel_scriptures_JST_NTR1.tex}
	
% ----------------------
\subsection{New Testament Revision 2}
% ----------------------
	\label{JST-NTR2}
	
	\input{gospel_scriptures_JST_NTR2.tex}